% Section 6: Discussion
\section{Discussion}
\label{sec:discussion}

\subsection{Key Findings}

\textbf{The LC correction preserves capability ordering.} $r_{\mathrm{partial}} = 0.985$
provides empirical validation of AlpacaEval 2.0's design goal across 223 diverse models
spanning GPT-4 class to lightweight instruction-tuned variants.

\textbf{Why partial exceeds bivariate.} Verbosity acts as a mild suppressor
($\rho(\texttt{win\_rate}, \texttt{avg\_length}) \approx 0.63$). Partialling out verbosity
reveals the capability signal more cleanly --- strengthening from 0.94 to 0.985. This is
consistent with \citet{Hu2024Explaining}: the desirability (length-independent) channel in
\texttt{win\_rate} survives LC correction without attenuation.

\textbf{The verbosity penalty.} $\beta_{\mathrm{len}} = -4.37$ (negative) confirms the LC
correction penalizes verbosity once capability is controlled --- exactly the intended behavior.

\subsection{The $\Delta$ Nuance}

Non-monotonic Q3 median $\Delta = 5.43 >$ Q4 $= 0.91$ reflects two factors:
(1) mathematical dependency ($\Delta = \text{LC} - \text{win}$; grouping by \texttt{win\_rate}
quartiles makes $\Delta$ partly self-referential); (2) possible Q3 verbosity-exploitation
heterogeneity. The $10\times$ effect-size gap between H-M3 ($\varepsilon^2 = 0.088$) and H-C1
($\varepsilon^2 = 0.883$) establishes \texttt{LC\_winrate} as the appropriate DV for
capability-alignment studies.

\subsection{Limitations}

\textbf{L1 --- Observational study}: Association, not causation. Unmeasured confounders (model
family, RLHF budget, base model) could co-vary with both metrics.

\textbf{L2 --- Single dataset}: AlpacaEval 2.0 (2023--2024 era models). Results may not
generalize to MT-Bench, MMLU, Chatbot Arena, or post-2024 frontier models.

\textbf{L3 --- Imperfect capability proxy}: \texttt{win\_rate} conflates capability with
response style preferences and annotator biases.

\textbf{L4 --- Mild heteroscedasticity}: BP $p = 0.011$ in H-M1. Non-critical given
$p_{\mathrm{win}} = 4.58 \times 10^{-145}$.

\textbf{L5 --- $\Delta$-monotonicity not confirmed}: Step 3 of original hypothesis partially
supported; \texttt{LC\_winrate}-based evidence (H-C1) provides stronger support.

\subsection{Broader Impact}

This work validates widely used LLM evaluation infrastructure. Positive impacts: more confident
use of LC-based rankings; FWL verification as replicable confound-control template; principled
DV choice guidance. No personal data used; all analyses on publicly available aggregated
leaderboard scores. No harmful applications apparent.
